\chapter{Background}

\section{LLM-based Agents}
An agent based on a large language model alternates between generating text and calling
tools. The model receives a task, a description of the tools it may call and the results of
previous calls, and decides the next step. Everything the agent knows about an external
system reaches it as text in this loop: the descriptions of the tools and what the tools
return. The amount of that text is the context the agent consumes, and it is billed and
limited. ReAct showed that interleaving reasoning steps with actions on external tools
improves question answering and decision making \cite{yao2022react}, and Toolformer showed
that a model can learn by itself when and how to call tools such as a search engine or a
calculator \cite{schick2023toolformer}.

\section{Model Context Protocol}
The Model Context Protocol (MCP) is an open standard for connecting AI applications to
external systems \cite{mcp2026intro}. A server can expose three kinds of primitives: tools,
which are executable functions the application can invoke; resources, which are data sources
that provide context; and prompts, which are reusable templates. The server studied here
exposes its catalogue through four tools and one resource. The protocol was created at
Anthropic and released as open source in November 2024 \cite{anthropic2024mcp}.

\section{Data Contracts}
A data contract is a document, kept with the data, that states what a dataset contains and
what its consumers may rely on. The Open Data Contract Standard (ODCS) defines the keys and
values of a contract written in YAML \cite{odcs2026}. Its sections cover fundamentals,
schema, data quality, service-level agreements, infrastructure and servers, team and roles,
and custom properties, among others. The standard is maintained by Bitol, a project of the
Linux Foundation AI \& Data \cite{bitol2026}. The contracts of the case declare version
3.1.0, released in December 2025; version 3.2.0 followed in September 2026
\cite{odcsreleases2026}.

The contracts of the case use the custom properties to carry what the standard does not
define: tags and retrieval context for agents, example questions, and a \texttt{semantics}
block with the row grain, the dimension columns, the metrics with their quantity kind, unit
and aggregation rules, and the aliases of coded values. Figure~\ref{fig:contract} shows an
excerpt.

\begin{figure}[ht]
\centering
\begin{minipage}{0.88\textwidth}
\footnotesize
\begin{verbatim}
- property: semantics
  value:
    family: carga
    temporal_perspective: verificado
    row_grain:
      temporal: 1 dia
      entity: subsistema
      one_row_per: [din_instante, id_subsistema]
    metrics:
    - id: cargaenergiamwmed
      label: Carga Verificada
      expression: TRY_CAST(val_cargaenergiamwmed AS DOUBLE)
      quantity_kind: average_power
      unit_native: MWmed
      interval_hours: 24.0
      aggregation_modes:
        total_energy:
          sql: SUM((...) * 24.0) / 1000
          unit: GWh
      aliases: [carga, demanda, consumo energia, ...]
    query_routing:
      supports_time_series: true
      supports_ranking: false
\end{verbatim}
\end{minipage}
\caption{Excerpt of the \texttt{semantics} block of the contract of the daily load dataset
(\texttt{carga-energia}), abridged.}
\label{fig:contract}
\end{figure}

\section{Ontologies}
An ontology is an explicit specification of the concepts of a domain and of the relations
between them \cite{gruber1993translation,studer1998knowledge}. In the RDF family of languages, the schema-level statements
(classes and properties) are called the T-Box and the statements about individuals the
A-Box. This work reuses five vocabularies. DCAT is an RDF vocabulary designed to facilitate
interoperability between data catalogues published on the Web, and is a W3C Recommendation
since August 2024 in its third version \cite{dcat3}. The RDF Data Cube vocabulary, a W3C
Recommendation of January 2014, describes multi-dimensional statistical data with measures
and dimensions \cite{datacube}. SKOS provides labels and concept schemes \cite{skos2009}, and
the CSVW metadata vocabulary describes tabular data and its columns \cite{csvw2015}. QUDT provides semantic specifications
for units of measure, quantity kinds, dimensions and data types \cite{qudt}.

For power systems, the Common Information Model (CIM) of the IEC defines classes such as
substation, generating unit, hydro power plant and reservoir \cite{iec61970301}. The
alignments of this work were checked against version 17 of the model, the one described by
that edition of the standard.
The ontology of this work aligns its spatial entities to CIM classes with loose matches and
keeps the concepts specific to the Brazilian system in its own namespace.
